Dans le cas ou l'on prend en compte le type de trajectoire il faut
moyenner les equations selon la relation \[
\left\{ \mathcal{A}\right\} =\frac{1}{\lambda\left(\xi_{0}\right)}\left[\frac{1}{2}\sum_{\sigma=\pm1}\right]_{T}\int_{-\theta_{c}}^{\theta_{c}}\frac{d\theta}{2\pi}\frac{\xi_{0}}{\xi}\mathcal{A}\]


ou \[
\lambda\left(\xi_{0}\right)=\int_{-\theta_{c}}^{\theta_{c}}\frac{d\theta}{2\pi}\frac{\xi_{0}}{\xi}\]
 et $\sigma=\frac{p_{\Vert}}{\left|p_{\Vert}\right|}$est le signe
de $p_{\Vert}$et $\theta$est l'angle poloidal.Ici $\theta_{c}$vaut
$\pi$pour les particules circulantes et $\theta_{T}=$ $\cos^{-1}\left(1-2\xi_{0}^{2}/\xi_{0T}^{2}\right)$pour
les particules piegees, avec $\xi_{0T}^{2}=$ $\frac{2\varepsilon}{1+\varepsilon},$
$\varepsilon=r/R_{0}$etant l'inverse du rapport d'aspect.L'equation
de Fokker-Planck moyennee devient alors\[
\left\{ \frac{\partial\left(rp^{2}f_{0}\right)}{\partial t}\right\} +\left\{ r\frac{\partial\left(p^{2}S_{p}\right)}{\partial p}\right\} -\left\{ rp\frac{\partial\left(\sqrt{1-\xi^{2}}S_{\xi}\right)}{\partial\xi}\right\} =0\]
 ou le parametre temps $t$evolue sur une echelle bien plus longue
que la periode de rebond \[
\tau_{t,b}=\frac{2\pi R_{0}q\left(r\right)}{\left|\xi_{0}\right|v}\lambda\left(\xi_{0}\right)\]
 On definit ici le facteur de securite $q\left(r\right)$de la surface
de flux magnetique situee au rayon $r$sur laquelle la moyenne est
effectuee et $R_{0}$le grand rayon du plasma.Toutes les valeurs sont
determinees au point ou le champ magnetique estminimum, la ou toutes
les particules passent, soit notamment $\xi_{0}$pour l'angle de diffusion
angulaire et la fonction de distribution $f_{0}^{\left(0\right)}$.On
a ainsi \[
\left\{ \frac{\partial\left(rp^{2}f_{0}\right)}{\partial t}\right\} =rp^{2}\frac{\partial}{\partial t}\left[\frac{1}{\lambda\left(\xi_{0}\right)}\left[\frac{1}{2}\sum_{\sigma=\pm1}\right]_{T}\int_{-\theta_{c}}^{\theta_{c}}\frac{d\theta}{2\pi}\frac{\xi_{0}}{\xi}f_{0}\right]\]
 et comme $f_{0}$est par definition symetrique en $\theta,$\[
\left\{ \frac{\partial\left(rp^{2}f_{0}\right)}{\partial t}\right\} =\frac{\partial\left(rp^{2}f_{0}^{(0)}\right)}{\partial t}\]


De maniere evidente \[
\left\{ r\frac{\partial\left(p^{2}S_{p}\right)}{\partial p}\right\} =r\frac{\partial\left(p^{2}\left\{ S_{p}\right\} \right)}{\partial p}\]
 et pour le second terme de droite il faut utiliser la relation $\frac{\partial}{\partial\xi}=$
$\frac{\xi}{\Psi\xi_{0}}\frac{\partial}{\partial\xi_{0}},$sachant
que $\xi=$ $\sigma\sqrt{1-\Psi\left(1-\xi_{0}^{2}\right)}$avec $\Psi\left(r,\theta\right)=\frac{1+\varepsilon}{1+\varepsilon\cos\theta}.$Il
vient alors \[
\left\{ rp\frac{\partial\left(\sqrt{1-\xi^{2}}S_{\xi}\right)}{\partial\xi}\right\} =\frac{rp}{\lambda\left(\xi_{0}\right)}\left[\frac{1}{2}\sum_{\sigma=\pm1}\right]_{T}\int_{-\theta_{c}}^{\theta_{c}}\frac{d\theta}{2\pi}\frac{\xi_{0}}{\xi}\frac{\partial\left(\sqrt{1-\xi^{2}}S_{\xi}\right)}{\partial\xi}\]
 soit apres le changement de variable \[
\left\{ rp\frac{\partial\left(\sqrt{1-\xi^{2}}S_{\xi}\right)}{\partial\xi}\right\} =\frac{rp}{\lambda\left(\xi_{0}\right)}\left[\frac{1}{2}\sum_{\sigma=\pm1}\right]_{T}\int_{-\theta_{c}}^{\theta_{c}}\frac{d\theta}{2\pi}\frac{1}{\Psi}\frac{\partial\left(\sqrt{1-\xi^{2}}S_{\xi}\right)}{\partial\xi_{0}}\]
 Afin de permuter la derivee en $\xi_{0}$et la somme sur $\sigma$pour
les particules piegees, il faut introduire $\sigma$, de sorte que
\begin{eqnarray*}
\left\{ rp\frac{\partial\left(\sqrt{1-\xi^{2}}S_{\xi}\right)}{\partial\xi}\right\}  & = & \frac{rp}{\lambda\left(\xi_{0}\right)}\frac{\partial\left(\left[\frac{1}{2}\sum_{\sigma=\pm1}\right]_{T}\int_{-\theta_{c}}^{\theta_{c}}\frac{d\theta}{2\pi}\frac{\sigma}{\Psi}\sqrt{1-\xi^{2}}S_{\xi}\right)}{\partial\left|\xi_{0}\right|}\\
 & = & \frac{rp}{\lambda\left(\xi_{0}\right)}\frac{\partial\left(\sigma\left[\frac{1}{2}\sum_{\sigma=\pm1}\right]_{T}\int_{-\theta_{c}}^{\theta_{c}}\frac{d\theta}{2\pi}\frac{\sigma}{\Psi}\sqrt{1-\xi^{2}}S_{\xi}\right)}{\partial\xi_{0}}\end{eqnarray*}
 et il est possible d'exprimer le terme dans la derivee partielle
$\frac{\partial}{\partial\xi_{0}}$sous la forme simple \[
\sigma\left[\frac{1}{2}\sum_{\sigma=\pm1}\right]_{T}\int_{-\theta_{c}}^{\theta_{c}}\frac{d\theta}{2\pi}\frac{\sigma}{\Psi}\sqrt{1-\xi^{2}}S_{\xi}=\lambda\left(\xi_{0}\right)\sqrt{1-\xi_{0}^{2}}\sigma\left\{ \frac{\sigma\xi}{\sqrt{\Psi}\xi_{0}}S_{\xi}\right\} \]
 ce qui donne \[
\left\{ rp\frac{\partial\left(\sqrt{1-\xi^{2}}S_{\xi}\right)}{\partial\xi}\right\} =\frac{rp}{\lambda\left(\xi_{0}\right)}\frac{\partial\left(\lambda\left(\xi_{0}\right)\sqrt{1-\xi_{0}^{2}}\sigma\left\{ \frac{\sigma\xi}{\sqrt{\Psi}\xi_{0}}S_{\xi}\right\} \right)}{\partial\xi_{0}}\]


En posant formellement les nouvelles definitions des flux moyennes
sousla forme de leur valeur ou le champ magnetique $B$est minimum,
\[
\mathbf{S}_{\mathbf{p}}^{(0)}=\left\{ \begin{array}{l}
S_{p}^{(0)}=\left\{ S_{p}\right\} \\
S_{\xi_{0}}^{(0)}=\sigma\left\{ \frac{\sigma\xi}{\sqrt{\Psi}\xi_{0}}S_{\xi}\right\} \end{array}\right.\]
 on peut garder une forme conservative simple pour l'equation de Fokker-Planck\[
\frac{\partial\left(rp^{2}f_{0}^{(0)}\right)}{\partial t}+r\frac{\partial}{\partial p}(p^{2}S_{p}^{(0)})-\frac{rp}{\lambda\left(\xi_{0}\right)}\frac{\partial}{\partial\xi_{0}}\left(\sqrt{1-\xi_{0}^{2}}\lambda\left(\xi_{0}\right)S_{\xi_{0}}^{(0)}\right)=0\]


Il est egalement possible de garder la forme originale pour l'expressiondes
flux, \begin{eqnarray*}
S_{p} & = & -D_{pp}\frac{\partial f_{0}}{\partial p}+D_{p\xi}\frac{\sqrt{1-\xi^{2}}}{p}\frac{\partial f_{0}}{\partial\xi}+F_{p}f_{0}\\
S_{\xi} & = & -D_{\xi p}\frac{\partial f_{0}}{\partial p}+D_{\xi\xi}\frac{\sqrt{1-\xi^{2}}}{p}\frac{\partial f_{0}}{\partial\xi}+F_{\xi}f_{0}\end{eqnarray*}


Les moyennes sur la trajectoire des flux definies ci-dessus donnent
alors\begin{eqnarray*}
\left\{ -D_{pp}\frac{\partial f_{0}}{\partial p}\right\}  & = & -\left\{ D_{pp}\right\} \frac{\partial f_{0}^{(0)}}{\partial p}\\
\left\{ D_{p\xi}\frac{\sqrt{1-\xi^{2}}}{p}\frac{\partial f_{0}}{\partial\xi}\right\}  & = & \frac{\sqrt{1-\xi_{0}^{2}}}{p}\left\{ \frac{\sigma\xi}{\sqrt{\Psi}\xi_{0}}D_{p\xi}\right\} \frac{\partial f_{0}^{(0)}}{\sigma\partial\xi_{0}}\\
\left\{ F_{p}f_{0}\right\}  & = & \left\{ F_{p}\right\} f_{0}^{(0)}\\
-\sigma\left\{ \frac{\sigma\xi}{\sqrt{\Psi}\xi_{0}}D_{\xi p}\frac{\partial f_{0}}{\partial p}\right\}  & = & -\sigma\left\{ \frac{\sigma\xi}{\sqrt{\Psi}\xi_{0}}D_{\xi p}\right\} \frac{\partial f_{0}^{(0)}}{\partial p}\\
-\sigma\left\{ \frac{\sigma\xi}{\sqrt{\Psi}\xi_{0}}D_{\xi\xi}\frac{\sqrt{1-\xi^{2}}}{p}\frac{\partial f_{0}}{\partial\xi}\right\}  & = & -\sigma\frac{\sqrt{1-\xi_{0}^{2}}}{p}\left\{ \frac{\xi^{2}}{\Psi\xi_{0}^{2}}D_{\xi\xi}\right\} \frac{\partial f_{0}^{(0)}}{\sigma\partial\xi_{0}}\\
\sigma\left\{ \frac{\sigma\xi}{\sqrt{\Psi}\xi_{0}}F_{\xi}f_{0}\right\}  & = & \sigma\left\{ \frac{\sigma\xi}{\sqrt{\Psi}\xi_{0}}F_{\xi}\right\} f_{0}^{(0)}\end{eqnarray*}
 de sorte que l'on peut reecrire \[
\mathbf{S}_{\mathbf{p}}^{(0)}=-\overline{\overline{D}}_{p}^{(0)}\cdot\nabla_{p,\xi_{0}}f_{0}^{(0)}+\overline{\overline{F}}_{p}^{(0)}f_{0}^{(0)}\]
 avec les elements de la matrice de diffusion \begin{eqnarray*}
D_{pp}^{(0)} & = & \left\{ D_{pp}\right\} \\
D_{p\xi}^{(0)} & = & \sigma\left\{ \frac{\sigma\xi}{\sqrt{\Psi}\xi_{0}}D_{p\xi}\right\} \\
D_{\xi p}^{(0)} & = & \sigma\left\{ \frac{\sigma\xi}{\sqrt{\Psi}\xi_{0}}D_{\xi p}\right\} \\
D_{\xi\xi}^{(0)} & = & \left\{ \frac{\xi^{2}}{\Psi\xi_{0}^{2}}D_{\xi\xi}\right\} \end{eqnarray*}
 et les termes de convection \begin{eqnarray*}
F_{p}^{(0)} & = & \left\{ F_{p}\right\} \\
F_{\xi}^{(0)} & = & \sigma\left\{ \frac{\sigma\xi}{\sqrt{\Psi}\xi_{0}}F_{\xi}\right\} \end{eqnarray*}


Le gradient s'exprime alors suivant \[
\nabla_{p,\xi_{0}}=\left\{ \begin{array}{l}
\nabla_{p}=\frac{\partial}{\partial p}\\
\nabla_{\xi_{0}}=-\frac{\sqrt{1-\xi_{0}^{2}}}{p}\frac{\partial}{\partial\xi_{0}}\end{array}\right.\]


Les composantes du flux moyenne sur la periode de rebond dans l'espace
desimpulsions valent donc \[
\mathbf{S}_{\mathbf{p}}^{(0)}=\left\{ \begin{array}{l}
S_{p}^{(0)}=-D_{pp}^{(0)}\frac{\partial f_{0}^{(0)}}{\partial p}+\frac{\sqrt{1-\xi_{0}^{2}}}{p}D_{p\xi}^{(0)}\frac{\partial f_{0}^{(0)}}{\partial\xi_{0}}+F_{p}^{(0)}f_{0}^{(0)}\\
S_{\xi_{0}}^{(0)}=-D_{\xi p}^{(0)}\frac{\partial f_{0}^{(0)}}{\partial p}+\frac{\sqrt{1-\xi_{0}^{2}}}{p}D_{\xi\xi}^{(0)}\frac{\partial f_{0}^{(0)}}{\partial\xi_{0}}+F_{\xi}^{(0)}f_{0}^{(0)}\end{array}\right.\]


L'equation de flux devient alors \begin{eqnarray*}
\frac{\partial\left(p^{2}S_{p}^{(0)}\right)}{\partial p} & = & \frac{\partial}{\partial p}\left(-p^{2}D_{pp}^{(0)}\frac{\partial f_{0}^{(0)}}{\partial p}+p^{2}F_{p}^{(0)}f_{0}^{(0)}\right)\\
 &  & +\sqrt{1-\xi_{0}^{2}}\frac{\partial}{\partial p}\left(pD_{p\xi}^{(0)}\right)\frac{\partial f_{0}^{(0)}}{\partial\xi_{0}}+\sqrt{1-\xi_{0}^{2}}pD_{p\xi}^{(0)}\frac{\partial^{2}f_{0}^{(0)}}{\partial p\partial\xi_{0}}\\
\frac{1}{\lambda\left(\xi_{0}\right)}\frac{\partial\left(\lambda\left(\xi_{0}\right)\sqrt{1-\xi_{0}^{2}}S_{\xi}^{(0)}\right)}{\partial\xi_{0}} & = & \frac{1}{\lambda\left(\xi_{0}\right)}\frac{\partial}{\partial\xi_{0}}\left(\lambda\left(\xi_{0}\right)D_{\xi\xi}^{(0)}\frac{1-\xi_{0}^{2}}{p}\frac{\partial f_{0}^{(0)}}{\partial\xi_{0}}+\lambda\left(\xi_{0}\right)\sqrt{1-\xi_{0}^{2}}F_{\xi}f_{0}^{(0)}\right)\\
 &  & -\frac{1}{\lambda\left(\xi_{0}\right)}\frac{\partial}{\partial\xi_{0}}\left(\lambda\left(\xi_{0}\right)\sqrt{1-\xi_{0}^{2}}D_{\xi p}^{(0)}\right)\frac{\partial f_{0}^{(0)}}{\partial p}-\sqrt{1-\xi_{0}^{2}}D_{\xi p}^{(0)}\frac{\partial^{2}f_{0}^{(0)}}{\partial\xi_{0}\partial p}\end{eqnarray*}


En discretisant, puis en les regroupant comme dans le cas non-moyenne,
onobtient \begin{eqnarray*}
 &  & r_{l+1/2}p_{i+1/2}^{2}\frac{f_{0,i+1/2,j+1/2}^{\left(0\right)(k+1)}-f_{0,i+1/2,j+1/2}^{\left(0\right)(k)}}{\Delta t}\\
 &  & +r_{l+1/2}\frac{-p_{i+1}^{2}D_{pp,i+1,j+1/2}^{\left(0\right)}\left.\frac{\partial f_{0}^{\left(0\right)}}{\partial p}\right|_{i+1,j+1/2}^{(k+1)}+p_{i+1}^{2}F_{p,i+1,j+1/2}^{\left(0\right)}f_{0,i+1,j+1/2}^{\left(0\right)(k+1)}}{\Delta p_{i+1/2}}\\
 &  & +r_{l+1/2}\frac{p_{i}^{2}D_{pp,i,j+1/2}^{\left(0\right)}\left.\frac{\partial f_{0}^{\left(0\right)}}{\partial p}\right|_{i,j+1/2}^{(k+1)}-p_{i}^{2}F_{p,i,j+1/2}^{\left(0\right)}f_{0,i,j+1/2}^{\left(0\right)(k+1)}}{\Delta p_{i+1/2}}\\
 &  & +r_{l+1/2}\sqrt{1-\xi_{0j+1/2}^{2}}p_{i+1/2}D_{p\xi,i+1/2,j+1/2}^{\left(0\right)}\left.\frac{\partial^{2}f_{0}^{\left(0\right)}}{\partial p\partial\xi_{0}}\right|_{i+1/2,j+1/2}^{(k+1)}\\
 &  & +r_{l+1/2}\sqrt{1-\xi_{0j+1/2}^{2}}\left.\frac{\partial\left(pD_{p\xi}^{\left(0\right)}\right)}{\partial p}\right|_{i+1/2,j+1/2}\left.\frac{\partial f_{0}^{\left(0\right)}}{\partial\xi_{0}}\right|_{i+1/2,j+1/2}^{(k+1)}\\
 &  & -r_{l+1/2}\frac{D_{\xi\xi,i+1/2,j+1}^{\left(0\right)}\left(1-\xi_{0j+1}^{2}\right)\left.\frac{\partial f_{0}^{\left(0\right)}}{\partial\xi_{0}}\right|_{i+1/2,j+1}^{(k+1)}+p_{i+1/2}\sqrt{1-\xi_{0j+1}^{2}}F_{\xi,i+1/2,j+1}^{\left(0\right)}f_{0,i+1/2,j+1}^{\left(0\right)(k+1)}}{\Delta\xi_{0j+1/2}}\frac{\lambda\left(\xi_{0j+1}\right)}{\lambda\left(\xi_{0j+1/2}\right)}\\
 &  & +r_{l+1/2}\frac{D_{\xi\xi,i+1/2,j}^{\left(0\right)}\left(1-\xi_{0j}^{2}\right)\left.\frac{\partial f_{0}^{\left(0\right)}}{\partial\xi_{0}}\right|_{i+1/2,j}^{(k+1)}+p_{i+1/2}\sqrt{1-\xi_{0j}^{2}}F_{\xi,i+1/2,j}^{\left(0\right)}f_{0,i+1/2,j}^{\left(0\right)(k+1)}}{\Delta\xi_{0j+1/2}}\frac{\lambda\left(\xi_{0j}\right)}{\lambda\left(\xi_{0j+1/2}\right)}\\
 &  & +r_{l+1/2}p_{i+1/2}\sqrt{1-\xi_{0j+1}^{2}}D_{\xi p,i+1/2,j+1/2}^{\left(0\right)}\left.\frac{\partial^{2}f_{0}^{\left(0\right)}}{\partial\xi_{0}\partial p}\right|_{i+1/2,j+1/2}^{(k+1)}\\
 &  & +r_{l+1/2}p_{i+1/2}\left.\frac{\partial}{\partial\xi_{0}}\left(\lambda\left(\xi_{0}\right)\sqrt{1-\xi_{0}^{2}}D_{\xi p}^{\left(0\right)}\right)\right|_{i+1/2,j+1/2}\left.\frac{\partial f_{0}^{\left(0\right)}}{\partial p}\right|_{i+1/2,j+1/2}^{(k+1)}\\
 & = & 0\end{eqnarray*}
 ou l'indice $(k)$indique que la variable est prise au temps $t=k\Delta t,$
$\Delta t$etant le pas de temps.Le terme de droite est pris au temps
$(k+1)$, ce qui correspond a un schema totalement implicite ou temps
inverse, inconditionnellement stable numeriquement, point qui sera
rediscute dans un chapitresuivant.On discretise egalement \begin{eqnarray*}
\left.\frac{\partial^{2}f_{0}^{\left(0\right)}}{\partial p\partial\xi_{0}}\right|_{i+1/2,j+1/2}^{(k+1)} & = & \left.\frac{\partial^{2}f_{0}^{\left(0\right)}}{\partial\xi_{0}\partial p}\right|_{i+1/2,j+1/2}^{(k+1)}\\
 & = & \frac{f_{0,i+3/2,j+3/2}^{\left(0\right)(k+1)}-f_{0,i+3/2,j-1/2}^{\left(0\right)(k+1)}-f_{0,i-1/2,j+3/2}^{\left(0\right)(k+1)}+f_{0,i-1/2,j-1/2}^{\left(0\right)(k+1)}}{\left(\Delta p_{i+1}+\Delta p_{i}\right)\left(\Delta\xi_{0j+1}+\Delta\xi_{0j}\right)}\end{eqnarray*}
 tandis que \begin{eqnarray*}
\left.\frac{\partial\left(pD_{p\xi}^{\left(0\right)}\right)}{\partial p}\right|_{i+1/2,j+1/2} & = & \frac{p_{i+1}D_{p\xi,i+1,j+1/2}^{\left(0\right)}-p_{i}D_{p\xi,i,j+1/2}^{\left(0\right)}}{\Delta p_{i+1/2}}\\
\left.\frac{\partial}{\partial\xi_{0}}\left(\lambda\left(\xi_{0}\right)\sqrt{1-\xi_{0}^{2}}D_{\xi p}^{\left(0\right)}\right)\right|_{i+1/2,j+1/2} & = & \frac{\lambda\left(\xi_{0j+1}\right)\sqrt{1-\xi_{0j+1}^{2}}D_{\xi p,i+1/2,j+1}^{\left(0\right)}-\lambda\left(\xi_{0j}\right)\sqrt{1-\xi_{0j}^{2}}D_{\xi p,i+1/2,j}^{\left(0\right)}}{\Delta\xi_{0j+1/2}\lambda\left(\xi_{0j+1/2}\right)}\end{eqnarray*}


L'equation de Fokker-Planck moyennee sur la periode de rebond devient
alors sous forme discrete \[
r_{l+1/2}p_{i+1/2}^{2}\frac{f_{0,i+1/2,j+1/2}^{\left(0\right)\left(k+1\right)}-f_{0,i+1/2,j+1/2}^{\left(0\right)\left(k\right)}}{\Delta t}+r_{l+1/2}\sum_{i^{\prime}=i-1}^{i^{\prime}=i+1}\sum_{j^{\prime}=j-1}^{j^{\prime}=j+1}\overline{\overline{M}}_{p,i^{\prime}+1/2,j^{\prime}+1/2}^{\left(0\right)}f_{0,i^{\prime}+1/2,j^{\prime}+1/2}^{\left(0\right)(k+1)}=0\]
 , les coefficients de la matrice valant \begin{eqnarray*}
\overline{\overline{M}}_{p,i+3/2,j+3/2}^{\left(0\right)} & = & \frac{p_{i+1/2}}{\Delta p_{i+1}+\Delta p_{i}}\frac{\sqrt{1-\xi_{0j+1/2}^{2}}}{\Delta\xi_{0j+1}+\Delta\xi_{0j}}D_{p\xi,i+1/2,j+1/2}^{\left(0\right)}\\
 &  & +\frac{p_{i+1/2}}{\Delta p_{i+1}+\Delta p_{i}}\frac{\sqrt{1-\xi_{0j+1/2}^{2}}}{\Delta\xi_{0j+1}+\Delta\xi_{0j}}D_{\xi p,i+1/2,j+1/2}^{\left(0\right)}\end{eqnarray*}
\begin{eqnarray*}
\overline{\overline{M}}_{p,i+1/2,j+3/2}^{\left(0\right)} & = & \frac{p_{i+1}}{\Delta p_{i+1/2}}\frac{\sqrt{1-\xi_{0j+1/2}^{2}}}{\Delta\xi_{0j+1}+\Delta\xi_{0j}}D_{p\xi,i+1,j+1/2}^{\left(0\right)}\\
 &  & -\frac{p_{i}}{\Delta p_{i+1/2}}\frac{\sqrt{1-\xi_{0j+1/2}^{2}}}{\Delta\xi_{0j+1}+\Delta\xi_{0j}}D_{p\xi,i,j+1/2}^{\left(0\right)}\\
 &  & -\frac{1-\xi_{0j+1}^{2}}{\Delta\xi_{0j+1/2}\Delta\xi_{0j+1}}\frac{\lambda\left(\xi_{0j+1}\right)}{\lambda\left(\xi_{0j+1/2}\right)}D_{\xi\xi,i+1/2,j+1}^{\left(0\right)}\\
 &  & -p_{i+1/2}\frac{\sqrt{1-\xi_{0j+1}^{2}}}{\Delta\xi_{0j+1/2}}\left(1-\delta_{\xi,i+1/2,j+1}^{\left(0\right)}\right)\frac{\lambda\left(\xi_{0j+1}\right)}{\lambda\left(\xi_{0j+1/2}\right)}F_{\xi,i+1/2,j+1}^{\left(0\right)}\end{eqnarray*}
\begin{eqnarray*}
\overline{\overline{M}}_{p,i-1/2,j+3/2}^{\left(0\right)} & = & -\frac{p_{i+1/2}}{\Delta p_{i+1}+\Delta p_{i}}\frac{\sqrt{1-\xi_{0j+1/2}^{2}}}{\Delta\xi_{0j+1}+\Delta\xi_{0j}}D_{p\xi,i+1/2,j+1/2}^{\left(0\right)}\\
 &  & -\frac{p_{i+1/2}}{\Delta p_{i+1}+\Delta p_{i}}\frac{\sqrt{1-\xi_{0j+1/2}^{2}}}{\Delta\xi_{0j+1}+\Delta\xi_{0j}}D_{\xi p,i+1/2,j+1/2}^{\left(0\right)}\end{eqnarray*}
\begin{eqnarray*}
\overline{\overline{M}}_{p,i+3/2,j+1/2}^{\left(0\right)} & = & -\frac{p_{i+1}^{2}}{\Delta p_{i+1}\Delta p_{i+1/2}}D_{pp,i+1,j+1/2}^{\left(0\right)}\\
 &  & +\frac{p_{i+1/2}}{\Delta p_{i+1}+\Delta p_{i}}\frac{\sqrt{1-\xi_{0j+1}^{2}}}{\Delta\xi_{0j+1/2}}\frac{\lambda\left(\xi_{0j+1}\right)}{\lambda\left(\xi_{0j+1/2}\right)}D_{\xi p,i+1/2,j+1}^{\left(0\right)}\\
 &  & -\frac{p_{i+1/2}}{\Delta p_{i+1}+\Delta p_{i}}\frac{\sqrt{1-\xi_{0j}^{2}}}{\Delta\xi_{0j+1/2}}\frac{\lambda\left(\xi_{0j}\right)}{\lambda\left(\xi_{0j+1/2}\right)}D_{\xi p,i+1/2,j}^{\left(0\right)}\\
 &  & +\frac{p_{i+1}^{2}}{\Delta p_{i+1/2}}\left(1-\delta_{p,i+1,j+1/2}^{\left(0\right)}\right)F_{p,i+1,j+1/2}^{\left(0\right)}\end{eqnarray*}
\begin{eqnarray*}
\overline{\overline{M}}_{p,i+1/2,j+1/2}^{\left(0\right)} & = & \frac{p_{i+1}^{2}}{\Delta p_{i+1/2}\Delta p_{i+1}}D_{pp,i+1,j+1/2}^{\left(0\right)}\\
 &  & +\frac{p_{i+1}^{2}}{\Delta p_{i+1/2}}F_{p,i+1,j+1/2}^{\left(0\right)}\delta_{p,i+1,j+1/2}\\
 &  & +\frac{p_{i}^{2}}{\Delta p_{i+1/2}\Delta p_{i}}D_{pp,i,j+1/2}^{\left(0\right)}\\
 &  & -\frac{p_{i}^{2}}{\Delta p_{i+1/2}}\left(1-\delta_{p,i,j+1/2}^{\left(0\right)}\right)F_{p,i,j+1/2}^{\left(0\right)}\\
 &  & +\frac{1-\xi_{0j+1}^{2}}{\Delta\xi_{0j+1/2}\Delta\xi_{0j+1}}\frac{\lambda\left(\xi_{0j+1}\right)}{\lambda\left(\xi_{0j+1/2}\right)}D_{\xi\xi,i+1/2,j+1}^{\left(0\right)}\\
 &  & -p_{i+1/2}\frac{\sqrt{1-\xi_{0j+1}^{2}}}{\Delta\xi_{0j+1/2}}\delta_{\xi,i+1/2,j+1}^{\left(0\right)}\frac{\lambda\left(\xi_{0j+1}\right)}{\lambda\left(\xi_{0j+1/2}\right)}F_{\xi,i+1/2,j+1}^{\left(0\right)}\\
 &  & +\frac{1-\xi_{0j}^{2}}{\Delta\xi_{0j}\Delta\xi_{0j+1/2}}\frac{\lambda\left(\xi_{0j}\right)}{\lambda\left(\xi_{0j+1/2}\right)}D_{\xi\xi,i+1/2,j}^{\left(0\right)}\\
 &  & +p_{i+1/2}\frac{\sqrt{1-\xi_{0j}^{2}}}{\Delta\xi_{0j+1/2}}\left(1-\delta_{\xi,i+1/2,j}^{\left(0\right)}\right)\frac{\lambda\left(\xi_{0j}\right)}{\lambda\left(\xi_{0j+1/2}\right)}F_{\xi,i+1/2,j}^{\left(0\right)}\end{eqnarray*}
\begin{eqnarray*}
\overline{\overline{M}}_{i-1/2,j+1/2}^{\left(0\right)} & = & -\frac{p_{i}^{2}}{\Delta p_{i+1/2}\Delta p_{i}}D_{pp,i,j+1/2}^{\left(0\right)}\\
 &  & -\frac{p_{i+1/2}}{\Delta p_{i+1}+\Delta p_{i}}\frac{\sqrt{1-\xi_{0j+1}^{2}}}{\Delta\xi_{0j+1/2}}\frac{\lambda\left(\xi_{0j+1}\right)}{\lambda\left(\xi_{0j+1/2}\right)}D_{\xi p,i+1/2,j+1}^{\left(0\right)}\\
 &  & +\frac{p_{i+1/2}}{\Delta p_{i+1}+\Delta p_{i}}\frac{\sqrt{1-\xi_{0j}^{2}}}{\Delta\xi_{0j+1/2}}\frac{\lambda\left(\xi_{0j}\right)}{\lambda\left(\xi_{0j+1/2}\right)}D_{\xi p,i+1/2,j}^{\left(0\right)}\\
 &  & -\frac{p_{i}^{2}}{\Delta p_{i+1/2}}\delta_{p,i,j+1/2}^{\left(0\right)}F_{p,i,j+1/2}^{\left(0\right)}\end{eqnarray*}
\begin{eqnarray*}
\overline{\overline{M}}_{p,i+3/2,j-1/2}^{\left(0\right)} & = & -\frac{p_{i+1/2}}{\Delta p_{i+1}+\Delta p_{i}}\frac{\sqrt{1-\xi_{0j+1/2}^{2}}}{\Delta\xi_{0j+1}+\Delta\xi_{0j}}D_{p\xi,i+1/2,j+1/2}^{\left(0\right)}\\
 &  & -\frac{p_{i+1/2}}{\Delta p_{i+1}+\Delta p_{i}}\frac{\sqrt{1-\xi_{0j+1/2}^{2}}}{\Delta\xi_{0j+1}+\Delta\xi_{0j}}D_{\xi p,i+1/2,j+1/2}^{\left(0\right)}\end{eqnarray*}
\begin{eqnarray*}
\overline{\overline{M}}_{p,i+1/2,j-1/2}^{\left(0\right)} & = & -\frac{p_{i+1}}{\Delta p_{i+1/2}}\frac{\sqrt{1-\xi_{0j+1/2}^{2}}}{\Delta\xi_{0j+1}+\Delta\xi_{0j}}D_{p\xi,i+1,j+1/2}^{\left(0\right)}\\
 &  & +\frac{p_{i}}{\Delta p_{i+1/2}}\frac{\sqrt{1-\xi_{0j+1/2}^{2}}}{\Delta\xi_{0j+1}+\Delta\xi_{0j}}D_{p\xi,i,j+1/2}^{\left(0\right)}\\
 &  & -\frac{1-\xi_{0j}^{2}}{\Delta\xi_{0j+1/2}\Delta\xi_{0j}}\frac{\lambda\left(\xi_{0j}\right)}{\lambda\left(\xi_{0j+1/2}\right)}D_{\xi\xi,i+1/2,j}^{\left(0\right)}\\
 &  & +p_{i+1/2}\frac{\sqrt{1-\xi_{0j}^{2}}}{\Delta\xi_{0j+1/2}}\delta_{\xi,i+1/2,j}^{\left(0\right)}\frac{\lambda\left(\xi_{0j}\right)}{\lambda\left(\xi_{0j+1/2}\right)}F_{\xi,i+1/2,j}^{\left(0\right)}\end{eqnarray*}
\begin{eqnarray*}
\overline{\overline{M}}_{p,i-1/2,j-1/2}^{\left(0\right)} & = & \frac{p_{i+1/2}}{\Delta p_{i+1}+\Delta p_{i}}\frac{\sqrt{1-\xi_{0j+1/2}^{2}}}{\Delta\xi_{0j+1}+\Delta\xi_{0j}}D_{p\xi,i+1/2,j+1/2}^{\left(0\right)}\\
 &  & +\frac{p_{i+1/2}}{\Delta p_{i+1}+\Delta p_{i}}\frac{\sqrt{1-\xi_{0j+1/2}^{2}}}{\Delta\xi_{0j+1}+\Delta\xi_{0j}}D_{\xi p,i+1/2,j+1/2}^{\left(0\right)}\end{eqnarray*}
 avec les coefficients de diffusion et de convection \[
\overline{\overline{D}}_{p}^{\left(0\right)}\rightarrow\left\{ \begin{array}{l}
D_{pp,i+1,j+1/2}^{\left(0\right)}\\
D_{pp,i,j+1/2}^{\left(0\right)}\\
D_{p\xi,i+1,j+1/2}^{\left(0\right)}\\
D_{\xi p,i+1/2,j+1}^{\left(0\right)}\\
D_{p\xi,i+1/2,j+1/2}^{\left(0\right)}\\
D_{\xi p,i+1/2,j+1/2}^{\left(0\right)}\\
D_{p\xi,i,j+1/2}^{\left(0\right)}\\
D_{\xi p,i+1/2,j}^{\left(0\right)}\\
D_{\xi\xi,i+1/2,j+1}^{\left(0\right)}\\
D_{\xi\xi,i+1/2,j}^{\left(0\right)}\end{array}\right.\]
 et \[
\overline{\overline{F}}_{p}^{\left(0\right)}\rightarrow\left\{ \begin{array}{l}
F_{p,i+1,j+1/2}^{\left(0\right)}\\
F_{p,i,j+1/2}^{\left(0\right)}\\
F_{\xi,i+1/2,j+1}^{\left(0\right)}\\
F_{\xi,i+1/2,j}^{\left(0\right)}\end{array}\right.\]


Les coefficients $\delta_{p}^{\left(0\right)}$sont maintenant definis
a partir de $\left.\frac{F_{p}^{\left(0\right)}}{D_{pp}^{\left(0\right)}}\right|_{effectif}$,
les termes $\delta_{\xi}^{\left(0\right)}$restant quant a eux identiques
a $\delta_{\xi},$car $\Delta\xi=\Delta\xi_{0}.$
